\chapter{Introduction}\label{ch:introduction}
\thispagestyle{pagestyle}

\section{Motivation and Context}
\label{sec:motivation}

The automotive industry is undergoing a fundamental architectural shift.
Traditional vehicles are increasingly giving way to software-defined vehicles
(SDVs) --- platforms in which electronic control units are consolidated onto
high-performance compute hardware and the vehicle's behaviour is governed by
software layers that can be updated over the air \cite{EclipseLeda2023}. This
transformation enables capabilities that were previously inconceivable: dynamic
reconfiguration of driving characteristics, continuous safety-critical updates
without a workshop visit, and the integration of machine-learning-based
components directly into the vehicle runtime.

At the same time, deep reinforcement learning (RL) has matured into a practical
engineering tool. Algorithms such as Proximal Policy Optimization (PPO)
\cite{Schulman2017} achieve human-competitive or superhuman performance in
continuous control tasks, with comparatively modest tuning effort. The natural
question that arises is: \textit{how can an RL-trained policy be moved,
end-to-end, from a training simulation onto an SDV runtime stack in a
reproducible and maintainable way?}

This question is not merely academic. Automotive-grade deployment requires that
a trained model run inside a standardised software container, interact with
vehicle signals through a well-defined data layer, and remain subject to
functional-safety monitors that can intervene when the model operates outside
its intended envelope. Bridging this gap between laboratory RL and production
SDV middleware is the core motivation behind the work presented in this thesis.

\section{Problem Statement}
\label{sec:problem}

Existing RL research in the autonomous driving domain concentrates almost
exclusively on training performance --- reward curves, sample efficiency, and
policy robustness \cite{Kiran2021}. The software engineering challenge of
taking a trained checkpoint and deploying it as a first-class citizen of an SDV
runtime stack is largely unexplored in the open-source community. Concretely,
the following sub-problems remain unaddressed in an integrated, open-source
pipeline:

\begin{enumerate}[leftmargin=2cm,topsep=1.15pt,itemsep=1.15pt,partopsep=1.15pt,parsep=1.15pt]
  \item \textbf{Simulation-to-middleware coupling.} A physics simulator and an
    SDV middleware stack speak fundamentally different protocols; a
    communication bridge that is both low-latency and decoupled from either
    system is needed.
  \item \textbf{Standardised vehicle-signal representation.} Raw simulator
    outputs must be mapped onto industry-standard Vehicle Signal Specification
    (VSS) paths so that the SDV safety layer can interpret them without
    simulator-specific knowledge.
  \item \textbf{Operational safety monitoring.} A deployed RL policy may
    extrapolate outside its training distribution; a lightweight, declarative
    safety monitor must be able to flag or override the policy's actions in
    real time.
  \item \textbf{Reproducible training workflow.} The entire pipeline --- from
    launching the simulation to saving a deployable checkpoint --- must be
    reproducible from a single control surface without manual orchestration.
\end{enumerate}

\section{Objectives}
\label{sec:objectives}

This thesis addresses the above problems through the design, implementation, and
evaluation of \textbf{SHILATE} (Software-defined-vehicle Hardware-In-the-Loop
Autonomous Training Environment). The specific objectives are:

\begin{enumerate}[leftmargin=2cm,topsep=1.15pt,itemsep=1.15pt,partopsep=1.15pt,parsep=1.15pt]
  \item Design and implement a Unity-based physics simulation of a rear-wheel-drive
    vehicle navigating an instrumented obstacle course.
  \item Define a shaped reward function that guides a PPO agent to complete laps
    reliably and safely, and expose a Gymnasium-compliant Python interface over MQTT.
  \item Implement a real-time health-monitoring and training-control layer inside
    the Unity Editor so that a researcher can observe and intervene in a training
    run without leaving the development environment.
  \item Construct an SDV deployment chain using Eclipse Leda, the Kuksa vehicle
    signal databroker, and a Velocitas vehicle application that enforces
    operational limits on the running policy.
  \item Evaluate the pipeline quantitatively: training convergence, inference
    lap performance, and end-to-end signal propagation through the SDV stack.
\end{enumerate}

\section{Thesis Structure}
\label{sec:structure}

The remainder of this document is organised as follows.
\Cref{ch:background} reviews theoretical background in reinforcement learning,
SDV middleware, and communication protocols, and situates SHILATE among related
work.
\Cref{ch:architecture} presents the overall system architecture.
\Cref{ch:unity} describes the Unity simulation in detail.
\Cref{ch:python} documents the Python RL training pipeline.
\Cref{ch:sdv} explains the SDV integration through Eclipse Leda.
\Cref{ch:experiments} reports experimental results.
\Cref{ch:conclusions} concludes with a summary of contributions and future work.


\section{\LaTeX\ Instructions}
\subsection{Structure} \label{section:structure}
Each file belonging to the document can be found in \texttt{bachelors\_en/}.
The main file of the template is \texttt{main.tex}, which must be compiled using pdflatex (usually the default option). Files in \texttt{components/} define the template and \textbf{should not be modified}.

In \texttt{customs.tex} reside the personal data which appear in places not directly accessible by the user and should be modified. Every file in \texttt{chapters/} represents a chapter where the content of the thesis is written. Under \texttt{images/} will be placed all the images used in the document.

\subsection{Authenticity Declaration}
The blank declaration can be obtained in two ways:
\begin{itemize}
	\item from \texttt{main.tex} residing in \texttt{authenticity\_declaration/}, an unsigned variant is created after filling personal data in \texttt{customs.tex}, similarly with the structure of \cref{section:structure}.
	\item from the document at \url{https://ac.upt.ro/finalizare-studii/}, in section \enquote{Documente importante}. This should be manually filled with personal data.
\end{itemize}
Regardless of the option chosen, the document must be printed, \textbf{physically signed(with a pen)} and scanned in order to be attached. In \texttt{customs.tex}, the path must be changed from the example to the signed and scanned PDF document, which will be saved in \texttt{bachelors\_ro/}.

\subsection{Other \LaTeX\ elements}
References to other elements in the document (excepting citations and links): \cref{section:introduction}; \cref{fig:myfig}.

For each reference, a \texttt{\textbackslash label} is needed, which will be inserted anywhere in the text or immediatley after chapter/section titles, and in figures/tables/code fragments.
\\[\baselineskip]

Citations will be added using bibtex, example: \cref{example:citation}.
Inserting images (\cref{fig:myfig}), tables (\cref{table:table1}) and code fragments (\cref{code:polym}) will be added using the corresponding code. For tables, an external formatting tool is recommended: (example \cref{example:table_url}).

\section{General Information}

Each chapter must have a clear structure, will begin on a new page and will have a title. It will be followed by two blank 12 pt. lines.

Each sub-section title (e.g. 1.2 General information) shall begin after a 12 points blank line after the text and shall be followed by a 12 points blank line.

The text of the paper will be justified. It is recommended to check the spelling of the text with the help of the speller of the Word program. It is recommended that the thesis paper should not exceed a number of 100 pages, annexes included.


Rules applying to the text of the paper:
\begin{enumerate}[leftmargin=2cm,topsep=1.15pt,itemsep=1.15pt,partopsep=1.15pt,parsep=1.15pt,label=\alph*.]
   \item Margins of the page – the following values will be used:
   \begin{itemize}[topsep=1.15pt,itemsep=1.15pt,partopsep=1.15pt,parsep=1.15pt]
     \item interior: 2 cm 
     \item exterior: 2 cm 
     \item up: 2,5 cm (header included)
     \item down: 2 cm
   \end{itemize}
   \item Line spacing – the text shall apply a 1.15 line spacing 
   \item Indentation – text within regular paragraphs will be aligned between the left and the right margins (justified). The first line of each paragraph will have a 1.5 cm indentation. There will be an exception for chapter titles, which will be aligned left, just like the titles of the tables and of the figures (see explanations below);
   \item Font – the 12 points Nimbus Sans font will be used, as well as the specific diacritics of the paper language (ex: ă, ş, ţ, î, â - for the Romanian language);
   \item Page numbering -  Page numbering runs from the title page to the last page of the paper, but the page number appears starts on the Introduction page. The page number is inserted at the bottom of the page, centered.
   \item f.	Page header – it shall be inserted starting with the introduction page and contains, in successive lines, the following text, aligned left and with the size of 8 points: (i) the text Politehnica University of Timișoara ; (ii) the name of the program of study and the year of the paper defense; (iii) the name of the candidate (left) and the title of the paper. At the right of the header, the UPT logo may be inserted; 
\end{enumerate}